\documentclass{leveldoc}
\DocUses{foundations,first_bend,slingshot,around_the_wall}
\DocFingerprint{2e4791b0671ac730}

\begin{document}

\section{What the level asks}\label{sec:asks}

The game's text:
\leveldescription

Three shots leave $(\val{xLaunch}, \val{yLaunch})$ with the same particle
($q = \val{q}$, $m = \val{m}$) and the same kinetic energy $T_0 = \val{Tzero}$.
The directions stored in the file are unit vectors at exactly $-12^\circ$, $0^\circ$
and $+12^\circ$. On every shot the detector is the rectangle
$\val{xDlo} \le x \le \val{xDhi}$, $\val{yDlo} \le y \le \val{yDhi}$.
An entry counts at every direction and every kinetic energy: the small height of the
box is the whole requirement. The player places at most $\val{maxCharges}$ charges,
of the particle's sign only, inside the region
$\val{xRegLo} \le x \le \val{xRegHi}$, $\val{yRegLo} \le y \le \val{yRegHi}$,
and no magnets. A flight ends when the particle touches a charge, leaves the world,
enters the detector or runs out of time
(\xsecref{foundations}{sec:board}); $t_{\max} = \val{tmax}$.

The reference solution adds $Q = \val{Qlo}$ at $(\val{xLo}, \val{yLo})$ and
$Q = \val{Qhi}$ at $(\val{xHi}, \val{yHi})$. Both have the particle's sign, so the
potential between them is positive and the rays slow down as they approach: a
decelerating lens, staggered and unequal. All three rays then enter the detector
through its left face. This document integrates those two-charge flights, and the
flight of each charge alone, by the fixed-step method of \levelref{around_the_wall}.

\section{Setup}\label{sec:setup}

The level file, field by field (the appendix reproduces it):
\setuptable

In symbols: $c = \val{c}$, the particle has radius $a = \val{aRad}$, and every charge
has radius $R = \val{R}$. The launch has
\[
  \gamma_0 = 1 + \frac{T_0}{mc^2} = \val{gammaZero}, \qquad
  \abs{\vect{p}_0} = \sqrt{\frac{T_0^2}{c^2} + 2mT_0} = \val{pZero}, \qquad
  v_0 = \frac{\abs{\vect{p}_0}}{\gamma_0 m} = \val{vZero}.
\]
The same $\gamma_0$, $\abs{\vect{p}_0}$ and $v_0$ start every shot; only the direction
changes. The world is $-1 \le x \le \val{xWorld}$, $-1 \le y \le \val{yWorld}$.
The two couplings are $\kappa_\mathrm{l} = qQ_\mathrm{l} = \val{kLo}$ and
$\kappa_\mathrm{h} = qQ_\mathrm{h} = \val{kHi}$. The level file gives the charges in
that order.

\begin{figure}[htbp]
\centering
\begin{tikzpicture}
\begin{axis}[axis equal image, width=0.88\linewidth, xmin=-1.8, xmax=32.4, ymin=-1.8, ymax=21.8,
  xtick={0,5,...,30}, ytick={0,5,...,20}, tick label style={font=\small},
  xlabel={$x$}, ylabel={$y$}, label style={font=\small}]
  \draw[black!45, dashed] (-1,-1) rectangle (31,21);
  \draw[black!50] (0,0) rectangle (30,20);
  \draw[black!40, dashed] (8,1) rectangle (18,19);
  \fill[green!50!black, fill opacity=0.16] (28,9) rectangle (30,11);
  \draw[green!40!black] (28,9) rectangle (30,11);
  \fill[red!70!black, fill opacity=0.35] (\rawval{xLo},\rawval{yLo}) circle[radius=0.3];
  \draw[red!70!black] (\rawval{xLo},\rawval{yLo}) circle[radius=0.3];
  \fill[red!70!black, fill opacity=0.35] (\rawval{xHi},\rawval{yHi}) circle[radius=0.3];
  \draw[red!70!black] (\rawval{xHi},\rawval{yHi}) circle[radius=0.3];
  \addplot[thick, orange!80!black] table {pathdown.dat};
  \addplot[very thick, blue!70!black] table {pathaxis.dat};
  \addplot[thick, purple!70!black] table {pathup.dat};
  \draw[-{Stealth[length=6pt]}, thick, orange!80!black] (\rawval{xLaunch},\rawval{yLaunch}) --
    (\rawval{arrowDx},\rawval{arrowDy});
  \draw[-{Stealth[length=6pt]}, very thick, blue!70!black] (\rawval{xLaunch},\rawval{yLaunch}) --
    (\rawval{arrowX},\rawval{arrowY});
  \draw[-{Stealth[length=6pt]}, thick, purple!70!black] (\rawval{xLaunch},\rawval{yLaunch}) --
    (\rawval{arrowUx},\rawval{arrowUy});
  \fill (\rawval{xLaunch},\rawval{yLaunch}) circle[radius=0.11];
  \node[font=\scriptsize, red!60!black, anchor=north] at (8,6.55) {$Q_\mathrm{l}$};
  \node[font=\scriptsize, red!60!black, anchor=south] at (12,15.45) {$Q_\mathrm{h}$};
  \node[font=\scriptsize, green!30!black, anchor=south] at (29,11.15) {detector};
\end{axis}
\end{tikzpicture}
\caption{The three rays of the reference. Orange is $-12^\circ$, blue is the axis, purple
is $+12^\circ$. All three enter the left face of the detector. The discs are the two
repelling charges, each of radius $R$. The dashed rectangle is the region in which the
player may place charges.}
\label{fig:board}
\end{figure}

\section{Principles and model}\label{sec:model}

\begin{principles}
  \item The field of a charged sphere and the force on a non-polarizable particle,
    \xeqref{foundations}{eq:coulomb} and \xeqref{foundations}{eq:sphereforce}.
  \item The relativistic equation of motion \xeqref{foundations}{eq:eom}, the launch
    \xeqref{foundations}{eq:launch} and the energy \xeqref{foundations}{eq:energy}. The
    motion stays in the plane (\xsecref{foundations}{sec:plane}).
  \item The board's rules (\xsecref{foundations}{sec:board}). The first entry into the
    detector decides. This level adds no direction and no kinetic window.
  \item The classical Runge--Kutta method \xeqref{foundations}{eq:rk4} and the fixed-step
    flight of \xsecref{around_the_wall}{sec:method}. With two centres the orbit equation
    does not separate (\xeqref{slingshot}{eq:noseparation}).
  \item The radiated power \xeqref{foundations}{eq:larmor}, with the acceleration
    \xeqref{first_bend}{eq:acceleration}, summed by Simpson's rule
    \xeqref{foundations}{eq:simpson}.
\end{principles}

The level has no magnets, so $\vect{B} = \vect{0}$. The reference potential is the sum
of the two spheres,
\begin{equation}\label{eq:potential}
  V(\vect{x})
  = \frac{\kappa_\mathrm{l}}{\abs{\vect{x} - \vect{r}_\mathrm{l}}}
  + \frac{\kappa_\mathrm{h}}{\abs{\vect{x} - \vect{r}_\mathrm{h}}}.
\end{equation}
The force is \xeqref{around_the_wall}{eq:force} with $N = \val{nCharges}$, the charges
in the file's order. The state $(x, y, p_x, p_y)$ obeys \xeqref{foundations}{eq:eom}
with that force until the first event. The energy is
\xeqref{around_the_wall}{eq:energy}, so the kinetic energy at a point is
\begin{equation}\label{eq:kinetic}
  T(\vect{x}) = W - mc^2 - V(\vect{x}).
\end{equation}

\begin{assumption}{The charges are held in place}\label{as:fixed}
None of the charges recoils. That is a rule of the game
(\xsecref{foundations}{sec:radiation}).
\end{assumption}

\begin{assumption}{No radiation}\label{as:radiation}
The radiated energy is not taken from the particle. Over the nine flights of this
document the largest fraction of $T_0$ is $\val{radFracMax}$ (§\ref{sec:hand}). The
game makes the same approximation and reports the same integral (\physref{8}).
\end{assumption}

\begin{assumption}{A point particle}\label{as:point}
$a = 0$, and the particle does not polarize, so \xeqref{foundations}{eq:sphereforce} is
exact. On the three reference rays the closest approach to either charge is
$\val{rClose}$, on the \val{rCloseRay} ray, to the \val{rCloseCharge} charge. The
clearance outside the sphere is $\val{clearClose}$. The samples that locate it are at
most $\val{sampleArc}$ apart (§\ref{sec:method}), far below that clearance.
\end{assumption}

\begin{assumption}{The step resolves the lens}\label{as:step}
The reported step is $h = \val{H}$. Halving it estimates the largest error of an end
state as $\val{errMax}$ (§\ref{sec:hand}). That is far below the height of the
detector, from $\val{yDlo}$ to $\val{yDhi}$. The tolerance used in
§\ref{sec:comparison} is the larger of $10^{-9}$ and ten times this estimate.
\end{assumption}

\section{Derivation}\label{sec:derivation}

\subsection{What the energy allows}

At the launch, \eqref{eq:potential} gives $V(\vect{x}_0) = \val{VZero}$, so
\[
  W = \val{W}, \qquad W - mc^2 = T_0 + V(\vect{x}_0) = \val{Wbind}.
\]
This is positive. Far from both charges $V$ falls to $0$, which lies below the bound,
so \eqref{eq:kinetic} allows the particle to leave the world. The integration has to
decide whether it does. On the axis ray the force at the launch is
$(\val{FxZero}, \val{FyZero})$. The horizontal part is negative: both charges stand
ahead of the particle and repel, so the lens retards the beam. The lower charge sits
below the launch and the upper charge above it, so their transverse pushes oppose, and
the sum is the small positive $\val{FyZero}$.

Both couplings are positive, so $V$ is positive everywhere and largest near the charges.
Along a ray that passes between them, \eqref{eq:kinetic} falls and then rises again.
That dip is the decelerating lens. The axis ray's minimum, computed in
§\ref{sec:hand}, is $\val{TminAxis}$, below $T_0$. At the detector the potential has
fallen below its value at the launch, and the arrival kinetic energies lie above $T_0$.

\subsection{Why the flight is integrated}\label{sec:noseparation}

With one charge, angular momentum closes the motion into the orbit of
\levelref{first_bend}. Two fixed centres already fail to separate in the relativistic
problem: the term \xeqref{slingshot}{eq:noseparation} does not split. The pair here is
staggered and unequal, so it has no symmetry that would remove the obstruction. Each
ray is the initial-value problem for \xeqref{around_the_wall}{eq:force} at
$N = \val{nCharges}$, and §\ref{sec:method} computes it to a stated accuracy.

\subsection{The method}\label{sec:method}

The state is advanced by the classical Runge--Kutta method \xeqref{foundations}{eq:rk4}
at the fixed step $h$, and the crossing of a boundary is the cubic Hermite interpolant
\xeqref{slingshot}{eq:hermite}, exactly as specified in
\xsecref{around_the_wall}{sec:method}. After each step the interpolant is sampled at
$\val{nSample}$ points. The speed is less than $c$, and the velocity terms of the
interpolant add less than $ch$, so the sample stays inside the box of the step's two
ends enlarged by $ch$. A sphere, the world's edge or the detector is tested only when
that box meets it. The samples are read in the board's order
(\xsecref{foundations}{sec:board}): a sphere of radius $R + a$ about a charge, then the
world's edge, then the detector. The earliest sample wins, and a tie keeps that order.
The first instant of the winning event is then the bisection of the interpolant,
$\val{nBisect}$ iterations. The position reported is the interpolant at that instant.
The momentum is one further step of the method, of that length.

When the event is a face of a box and the free coordinate lies in that face, the face
coordinate is set to the face value. The residual of the interpolant, before that
assignment, is the check that the step really did cross the face. The error of the
reported step is a fifteenth of its difference from the step $2h$
(\xsecref{foundations}{sec:numerics}). The energy
\xeqref{around_the_wall}{eq:energy} is a second check, one that does not use the game.

On each step the same interpolant is also sampled at nine equally spaced fractions,
including the endpoints. \eqref{eq:kinetic} at those samples gives the kinetic-energy
dip, and the distances to the two centres give the closest approach. The samples are an
eighth of a step apart in time. Since the speed is less than $c$, a point between them
lies within $\val{sampleArc}$ of a sample. The dip and the closest approach are
therefore located to that tolerance, which is coarser than the end-state estimate and
is used only for those two quantities.

\section{The reference solution by hand}\label{sec:hand}

\begin{procedure}
  \item\label{st:launch} \emph{The launch} (\xeqref{foundations}{eq:launch}):
    $\gamma_0 = \val{gammaZero}$ and $\abs{\vect{p}_0} = \val{pZero}$ on every shot,
    aimed at $-12^\circ$, along $+x$, and at $+12^\circ$. On the axis ray the force is
    $(\val{FxZero}, \val{FyZero})$, so the beam starts under a retarding force.
  \item\label{st:energy} \emph{The energy} \eqref{eq:kinetic}: $W = \val{W}$ and
    $W - mc^2 = \val{Wbind} > 0$. The world's edge is energetically open.
  \item\label{st:flight} \emph{The axis flight}: the method of §\ref{sec:method} with
    $h = \val{H}$. The states every five time units, and the drift of the energy:
    \begin{center}\small\val{states}\end{center}
    The kinetic energy falls to $\val{TminAxis}$ at $t = \val{tTmin}$, at
    $(\val{xTmin}, \val{yTmin})$. That point is $\val{rDipLo}$ from the lower charge
    and $\val{rDipHi}$ from the upper charge, between them. The lowest kinetic energies
    on the outer rays are $\val{RefDownTmin}$ and $\val{RefUpTmin}$.
  \item\label{st:entry} \emph{The three entries.} The axis step from
    $t_n = \val{tBefore}$, at $(\val{xBefore}, \val{yBefore})$, ends at
    $(\val{xAfter}, \val{yAfter})$. The interpolant crosses the \val{faceE} face at the
    fraction $u = \val{uE}$. The residual before the face coordinate is assigned is
    $\val{faceRes}$, so $x_\mathrm{e} = \val{xE}$, $y_\mathrm{e} = \val{yE}$ and
    $t_\mathrm{e} = \val{tE}$, after $\val{steps}$ steps. One step of length
    $t_\mathrm{e} - t_n$ gives $\vect{p}_\mathrm{e} = (\val{pxE}, \val{pyE})$, at
    $\val{angE}^\circ$ to $+x$.

    The $-12^\circ$ ray enters the \val{RefDownFace} face at
    $(\val{RefDownX}, \val{RefDownY})$ at $t = \val{RefDownT}$, with momentum
    $(\val{RefDownPx}, \val{RefDownPy})$, at $\val{RefDownAng}^\circ$.

    The $+12^\circ$ ray enters the \val{RefUpFace} face at
    $(\val{RefUpX}, \val{RefUpY})$ at $t = \val{RefUpT}$, with momentum
    $(\val{RefUpPx}, \val{RefUpPy})$, at $\val{RefUpAng}^\circ$.

    All three ordinates lie between $\val{yDlo}$ and $\val{yDhi}$.
  \item\label{st:check} \emph{The checks}: along the axis flight the energy stays within
    $\val{wDev}$ of $W$. The arrival kinetic energies from \eqref{eq:kinetic} are
    $\val{TarrDown}$, $\val{TarrAxis}$ and $\val{TarrUp}$, each above $\val{Tzero}$.
    Halving the step estimates the three end states as $\val{RefDownErr}$,
    $\val{errAxis}$ and $\val{RefUpErr}$. The same outcomes appear at $h$ and at $2h$.
  \item\label{st:rad} \emph{The radiation}: the power \xeqref{foundations}{eq:larmor}
    with the acceleration \xeqref{first_bend}{eq:acceleration}, summed by Simpson's rule
    on the steps of $h/2$. On the axis ray the integral is $\val{Erad}$, the fraction
    $\val{radFrac}$ of $T_0$. The relative difference from the same sum at step $h$,
    divided by fifteen, is $\val{radErr}$. The outer rays radiate $\val{RefDownRad}$
    and $\val{RefUpRad}$.
\end{procedure}

\begin{result}[The three rays]
All three enter the detector through its left face. The $-12^\circ$ ray arrives at
$(\val{RefDownX}, \val{RefDownY})$ at $t = \val{RefDownT}$, the axis ray at
$(\val{xE}, \val{yE})$ at $t = \val{tE}$, and the $+12^\circ$ ray at
$(\val{RefUpX}, \val{RefUpY})$ at $t = \val{RefUpT}$. Each ordinate lies between
$\val{yDlo}$ and $\val{yDhi}$. The level is solved.
\end{result}

By hand, Step~\ref{st:flight} is $\val{steps}$ steps of four evaluations of the force,
for the axis ray, and the same method for the two outer rays. The three step-halving
estimates are the ones in Step~\ref{st:check}. Across all nine flights of this document
the largest is $\val{errMax}$.

\section{Comparison with the game}\label{sec:comparison}

\comparisontable

The nine outcomes agree, and so do the nine end positions and times. The reference
momenta are compared as well. The tolerance on a state is $\val{tolFlight}$. Ten times
the largest step-halving estimate $\val{errMax}$ lies below $10^{-9}$, so the tolerance
is that floor. The three reference radiated energies are compared relatively, with
tolerance $10^{-4}$. The state tolerance and the radiation tolerance were fixed from
these estimates before the comparison.

The largest absolute difference is the time of the reference $+12^\circ$ ray,
$\val{cmpRefUpTDiff}$. Two ordinates of the upper charge alone equal it: the axis ray,
$\val{cmpHighAxisYDiff}$, and the $+12^\circ$ ray, $\val{cmpHighUpYDiff}$. The
difference is a few times $\val{errMax}$. The game's verified end states are good to
about $10^{-10}$ (\xsecref{foundations}{sec:verification}), so a difference of this
size is what that accuracy allows when two centres are summed.

The abscissae of the three reference entries agree exactly, as do the lower charge's
axis abscissa and the upper charge's downward and upward abscissae. The other three
abscissae differ by $\val{cmpLowDownXDiff}$, $\val{cmpLowUpXDiff}$ and
$\val{cmpHighAxisXDiff}$.

The axis ray's quadrature uncertainty is $\val{radErr}$, and the reference outer rays
give $\val{RefDownRadErr}$ and $\val{RefUpRadErr}$, all inside the relative tolerance.
The largest relative difference is the reference $-12^\circ$ ray,
$\val{cmpRefDownRadDiff}$.

\section{Other solutions}\label{sec:others}

Either added charge, without the other, fails to bring all three rays into the
detector. The method of §\ref{sec:method} integrates both (\code{placements.json}),
each ray at $h = \val{H}$ and at $2h$.
\begin{enumerate}[label=\arabic*.]
  \item $Q = \val{Qlo}$ at $(\val{xLo}, \val{yLo})$ only. All three rays leave through
    the right face of the world, above the detector.

    The $-12^\circ$ ray leaves at $(\val{LowDownX}, \val{LowDownY})$, at
    $t = \val{LowDownT}$.

    The momentum is $(\val{LowDownPx}, \val{LowDownPy})$, at $\val{LowDownAng}^\circ$
    to $+x$. The ordinate lies above $\val{yDhi}$. The step-halving estimate is
    $\val{LowDownErr}$. The radiated energy is $\val{LowDownRad}$.

    The axis ray leaves at $(\val{LowAxisX}, \val{LowAxisY})$, at $t = \val{LowAxisT}$.

    The momentum is $(\val{LowAxisPx}, \val{LowAxisPy})$, at $\val{LowAxisAng}^\circ$.
    The estimate is $\val{LowAxisErr}$. The radiated energy is $\val{LowAxisRad}$.

    The $+12^\circ$ ray leaves at $(\val{LowUpX}, \val{LowUpY})$, at
    $t = \val{LowUpT}$.

    The momentum is $(\val{LowUpPx}, \val{LowUpPy})$, at $\val{LowUpAng}^\circ$. The
    estimate is $\val{LowUpErr}$. The radiated energy is $\val{LowUpRad}$.
  \item $Q = \val{Qhi}$ at $(\val{xHi}, \val{yHi})$ only. The $-12^\circ$ ray and the
    axis ray leave through the right face of the world, below the detector. The
    $+12^\circ$ ray enters the detector.

    The $-12^\circ$ ray leaves at $(\val{HighDownX}, \val{HighDownY})$, at
    $t = \val{HighDownT}$.

    The momentum is $(\val{HighDownPx}, \val{HighDownPy})$, at
    $\val{HighDownAng}^\circ$. The ordinate lies below $\val{yDlo}$. The estimate is
    $\val{HighDownErr}$. The radiated energy is $\val{HighDownRad}$.

    The axis ray leaves at $(\val{HighAxisX}, \val{HighAxisY})$, at
    $t = \val{HighAxisT}$.

    The momentum is $(\val{HighAxisPx}, \val{HighAxisPy})$, at
    $\val{HighAxisAng}^\circ$. The estimate is $\val{HighAxisErr}$. The radiated energy
    is $\val{HighAxisRad}$.

    The $+12^\circ$ ray enters the \val{HighUpFace} face of the detector at
    $(\val{HighUpX}, \val{HighUpY})$, at $t = \val{HighUpT}$.

    The momentum is $(\val{HighUpPx}, \val{HighUpPy})$, at $\val{HighUpAng}^\circ$, so
    the ray is still heading downward as it enters. The estimate is $\val{HighUpErr}$.
    The radiated energy is $\val{HighUpRad}$.
\end{enumerate}

The reference needs both charges. The lower one turns the rays upward; alone, they
leave above the detector. The upper one turns them downward; alone, the $-12^\circ$ ray
and the axis ray leave below the detector, and the $+12^\circ$ ray enters while still
heading down. Together the three rays of Figure~\ref{fig:board} meet the box.
Section~\ref{sec:comparison} records the outcomes.

\section{What the game computes}\label{sec:engine}

The game integrates \xeqref{foundations}{eq:eom} with its eighth-order Runge--Kutta
method at the tolerances $10^{-10}$ and $10^{-12}$ (\physref{5}). It finds the crossing
on its interpolant (\physref{6}) and takes the first entry into the box as the outcome
(\xsecref{foundations}{sec:board}). It verifies that outcome by comparing its two runs
(\physref{7}). The model is this document's, including \ref{as:radiation} (\physref{8}).
The agreement in §\ref{sec:comparison} is between the fixed-step flight of
§\ref{sec:method} and that integration.

\section{Further reading}\label{sec:reading}

\begin{itemize}
  \item \levelref{around_the_wall}: the force of several centres and the fixed-step
    method with the enlarged-box test, used here at $N = \val{nCharges}$.
  \item \levelref{slingshot}: why two fixed centres have no closed relativistic orbit,
    and the Hermite crossing.
  \item \levelref{first_bend}: the acceleration that enters the radiated power.
  \item \xsecref{foundations}{sec:board} and \xsecref{foundations}{sec:numerics}: the
    order of events, and the step-halving estimate.
\end{itemize}

\levelfileappendix

\end{document}
